\chapter{Coupling and Constructive Methods: \mstar, Push-and-Swap, Push-and-Rotate}
\label{ch:ch11}
% Function names for pseudocode are prefixed with "Mstar" or "Pas" (push-and-swap)
% to stay unique across the book.
\SetKwFunction{MstarJointAStar}{JointAStar}
\SetKwFunction{MstarIndependenceDetection}{IndependenceDetection}
\SetKwFunction{MstarSearch}{MStar}
\SetKwFunction{MstarBackprop}{Backpropagate}
\SetKwFunction{MstarNeighbours}{LimitedNeighbours}
\SetKwFunction{MstarCollide}{Collisions}
\SetKwFunction{MstarPath}{PathFromParents}
\SetKwFunction{PasPlan}{PushAndSwap}
\SetKwFunction{PasPush}{Push}
\SetKwFunction{PasSwap}{Swap}
\SetKwFunction{PasRotate}{Rotate}

\begin{objectives}
\begin{itemize}
  \item Describe the joint state space of $k$ agents, compute its size and its branching factor, and explain why a plain joint \astar is hopeless beyond a handful of agents.
  \item Explain operator decomposition and independence detection, and prove that independence detection with an optimal group solver returns an optimal plan.
  \item Trace \mstar by hand on a small grid: follow the individual policies, detect a collision, build and backpropagate the collision sets, and expand only the agents that need it.
  \item State the conditions under which \mstar is complete and optimal, and sketch why.
  \item Describe the push, swap and rotate primitives, state what Push-and-Swap and Push-and-Rotate guarantee, and explain why their plans can be far from optimal.
  \item Place every multi-agent path finding method of the book in one comparison table and choose the right one for a given swarm scenario.
\end{itemize}
\end{objectives}

% ---------------------------------------------------------------------
\section{Why this matters for a drone swarm}
\label{sec:ch11-motivation}

Picture twelve inspection drones leaving a hangar through a grid of aisles.
Ten of them have routes that never come near each other. Two of them fly
towards each other in the same aisle, and one of them will have to slip into
a bay to let the other pass. Every multi-agent path finding (MAPF) solver has
to answer the same question for this scene: which agents must be planned
\emph{together}? Prioritized planning (\cref{ch:ch08}) answers ``none'', plans
the drones one after another and hopes that the fixed order works. A search in
the \emph{joint} state space, where one node is the position of all twelve
drones at once, answers ``all of them''; it is complete and optimal, but this
chapter shows that its state space has more configurations than there are
atoms in the hangar. Conflict-based search (\cbs, \cref{ch:ch09}) answers
``nobody, until a conflict appears, and then only through constraints'', which
is why the global planner of the hybrid architecture in \cref{ch:ch24} is \cbs
or its bounded-suboptimal relative \ecbs (\cref{ch:ch10}).

This chapter completes the picture with the two remaining answers. The first
is to \textbf{couple} agents explicitly and only where needed:
Standley's independence detection solves groups of agents separately and
merges two groups only when their plans conflict, and Wagner and Choset's
\mstar searches the joint space but lets every agent follow its own shortest
path until a collision proves that a subset of agents must be coordinated.
For the ten independent drones of the hangar this costs almost nothing; for
the two in the aisle it costs a small two-agent search. The second answer
gives up optimality altogether. Push-and-Swap and Push-and-Rotate never search
the joint space: they move the agents one at a time with a handful of
\emph{primitives} and guarantee a solution whenever one exists and at least
two cells are empty. Their plans can be several times longer than necessary,
but they arrive in polynomial time even on dense, puzzle-like instances that
stop every optimal solver.

In the training plan (\cref{ch:appA}) \mstar has priority \priority{Medium}
and Push-and-Swap priority \priority{Awareness}: you should be able to explain
what they do, trace \mstar on a small instance, and above all know when
coupling is worth its price and when a feasible plan is all a swarm needs.
The chapter is also the natural place to put every MAPF method of
\cref{ch:ch08,ch:ch09,ch:ch10} and of this chapter side by side, which
\cref{tab:ch11-comparison} does.

\paragraph{Roadmap.} \Cref{sec:ch11-intuition} gives the ideas.
\Cref{sec:ch11-problem} defines the joint state space, its size, the individual
policies and the collision sets. \Cref{sec:ch11-algorithm} presents joint
\astar with operator decomposition, independence detection and \mstar.
\Cref{sec:ch11-example} traces \mstar on a three-drone instance and compares
the effort of the four searches. \Cref{sec:ch11-properties} proves termination
and optimality. \Cref{sec:ch11-constructive} covers the push, swap and rotate
primitives and the completeness results of the constructive methods.
\Cref{sec:ch11-variants} lists the variants and contains the comparison table,
\cref{sec:ch11-implementation} the implementation notes, and
\cref{sec:ch11-drone} the place of these methods in the drone system.

% ---------------------------------------------------------------------
\section{The idea in plain words}
\label{sec:ch11-intuition}

Every agent knows its own shortest path. If the agents never met, the optimal
joint plan would simply be the bundle of these paths, and a search would not
be needed at all. Now imagine the joint state space of two agents as a
lattice: one axis measures how far agent~1 has progressed along its own
shortest path, the other how far agent~2 has (\cref{fig:ch11-idea}). The
bundle of individual paths is the diagonal of this lattice, a
\emph{one-dimensional} object inside a two-dimensional space. Joint \astar
ignores this structure and is prepared to expand any of the $|V|^2$
configurations. \mstar starts on the diagonal and stays there: from every
configuration it generates a single successor, the one in which every agent
takes its next policy step. Only when this successor is a collision does it
react. It marks the agents involved as a \emph{collision set}, sends the
mark back along the path it came from, and re-expands those configurations
with all combinations of moves of the marked agents, while the other agents
keep following their policies. The search widens into a small
two-dimensional patch, finds a detour, and once the agents have passed each
other it becomes one-dimensional again. Wagner and Choset call this
\textbf{subdimensional expansion}\index{subdimensional expansion}: the
dimension of the search space is raised locally, and only where a collision
has proved that it must be.

\begin{figure}[tb]
  \centering
  \inputfigure{ch11/idea}
  \caption{Subdimensional expansion for two agents. The grey dots are the
  $|V|^2$ joint configurations; the diagonal is the path in which both
  agents follow their individual policies, and \mstar expands nothing else
  until the policy successor is a collision (red cross). The collision set
  $\{1,2\}$ is backpropagated along the diagonal, the affected
  configurations are re-expanded with all joint moves of both agents (the
  blue patch), the detour in which agent~2 waits one step is found, and
  after the patch the search follows the diagonal again.}
  \label{fig:ch11-idea}
\end{figure}

Independence detection applies the same thought one level up. It solves
every agent alone, checks the resulting plans for conflicts, and merges two
groups into one joint search only if their plans actually conflict. Where
\mstar couples agents configuration by configuration, independence detection
couples them instance by instance. Both are exact: they return the same
optimal cost as the full joint search, because the agents they leave
decoupled are exactly those whose individual plans already fit together.

The constructive methods take a different stance. Push-and-Swap moves the
agents one at a time, like pebbles on the vertices of the graph. An agent
\emph{pushes} along its shortest path and shoves blocking agents into empty
vertices; when it meets an agent that has already reached its goal and cannot
be shoved, the two \emph{swap} places at a vertex with three neighbours.
Push-and-Rotate adds a \emph{rotate} primitive for fully occupied cycles.
These rules never look at the cost of the plan; their value is that they
always terminate, and that they find a plan whenever one exists on a graph
with two empty vertices.

\begin{keyidea}
Couple agents only where a collision proves that you must. Independence
detection merges groups whose optimal plans conflict; \mstar follows the
individual optimal policies through the joint space and branches only over
the agents in the collision set of a node. Both keep the guarantees of the
full joint search. If you need any plan rather than the best one, the push,
swap and rotate primitives build a valid plan in polynomial time, at a cost
that can be far from optimal.
\end{keyidea}

% ---------------------------------------------------------------------
\section{Problem statement and notation}
\label{sec:ch11-problem}

\subsection{The instance}

We use the classical MAPF setting of \cref{ch:ch07}: an undirected graph
$G = (V, E)$, here a 4-connected grid of free cells, and $k$
agents\index{agent} with distinct start vertices $s_1, \dots, s_k$ and
distinct goal vertices $\gamma_1, \dots, \gamma_k$. Time is discrete. In one
time step an agent either moves along an edge or waits in its cell. Two agents
are in a \textbf{vertex conflict}\index{conflict!vertex} if they occupy the
same cell at the same time and in a \textbf{swap conflict}\index{conflict!swap}
(edge conflict) if they exchange cells along an edge in the same time step.
Following an agent into the cell it is leaving is allowed, as in
\cref{ch:ch07}. A plan is a set of paths, one per agent, without conflicts;
an agent that has reached its goal stays there. Agents are numbered
$1, \dots, k$ in the text and $0, \dots, k-1$ in the code.

\subsection{The joint state space}

\begin{definition}[Joint configuration and joint graph]
\label{def:ch11-joint}
A \textbf{joint configuration}\index{joint configuration} is a tuple
$v = (v^1, \dots, v^k) \in V^k$ that gives the cell $v^i$ of every agent $i$.
It is \emph{collision-free} if $v^i \ne v^j$ for $i \ne j$. A \textbf{joint
move} from $v$ to $v'$ moves every agent to $v'^i \in \Succ(v^i) \cup
\set{v^i}$ (a neighbour or the same cell). The move is legal if $v'$ is
collision-free and no pair of agents swaps, $\neg(v'^i = v^j \wedge v'^j =
v^i)$ for all $i \ne j$. The \textbf{joint graph}\index{joint state space}
$G^k$ has the collision-free configurations as vertices and the legal joint
moves as edges. The joint start is $v_s = (s_1, \dots, s_k)$ and the joint
goal $v_f = (\gamma_1, \dots, \gamma_k)$. A path in $G^k$ from $v_s$ to $v_f$
is exactly a conflict-free plan, read one time step per edge.
\end{definition}

\begin{definition}[Cost of a joint move]
\label{def:ch11-cost}
A joint move from $v$ to $v'$ costs
\begin{equation}
  c(v, v') = \sum_{i=1}^{k} c_i(v^i, v'^i), \qquad
  c_i(u, u') = \begin{cases} 0 & \text{if } u = u' = \gamma_i,\\ 1 & \text{otherwise.}\end{cases}
  \label{eq:ch11-cost}
\end{equation}
Every move and every wait away from the goal costs one; resting at the goal
is free. The cost of a joint path is the sum of its move costs.
\index{cost!joint move}
\end{definition}

This is the convention of Standley~\cite{standley2010finding} and of Wagner
and Choset~\cite{wagner2011mstar}. It agrees with the sum of costs of
\cref{ch:ch07} on every plan in which no agent leaves its goal after arriving
there, which includes every plan of this chapter. It differs when an agent
must step off its goal to let another one through: the sum of costs then
also charges the free waits before the departure, so \cref{eq:ch11-cost} is
in general a lower bound on the sum of costs. The reason for adopting it
anyway is that it makes the cost of a move depend only on the two
configurations, which is what \astar needs (\cref{sec:ch11-implementation}
returns to this point).

\begin{figure}[tb]
  \centering
  \inputfigure{ch11/joint-graph}
  \caption{The joint graph of two agents on a four-vertex graph. Of the
  $4^2 = 16$ configurations, the four on the diagonal put both agents on the
  same vertex and are not states at all; the three dotted edges are joint
  moves in which the agents would swap along an edge; what remains are 12
  states and 18 joint moves (a wait of both agents is a self-loop and is not
  drawn). The optimal plan from $ac$ to $ca$ has cost 7: agent~2 steps aside
  into $d$, agent~1 waits once and passes, and agent~2 comes back.}
  \label{fig:ch11-joint-graph}
\end{figure}

\Cref{fig:ch11-joint-graph} draws the joint graph of two agents on a graph
with four vertices, generated by \code{gen\_ch11\_joint\_graph.py}. Even here
the structure is visible: vertex collisions delete states, swap conflicts
delete edges, and the plan is a path through what is left. The trouble is
the size. On a grid with $|V|$ free cells and $k$ agents there are $|V|^k$
configurations, of which $|V|(|V|-1)\cdots(|V|-k+1) = |V|^{\underline{k}}$
are collision-free, and a node of the joint graph has up to $(b+1)^k$
successors, where $b$ is the number of neighbours of a cell: $5^k$ on a
4-connected grid with the wait action. \Cref{tab:ch11-sizes} evaluates these
numbers for the $8 \times 8$ maps of the experiment in
\cref{sec:ch11-example}, which have about $54$ free cells, and for a
$20 \times 20$ map. Ten agents on the small map already have more joint
configurations than the number of seconds since the Big Bang, and every
expansion of joint \astar has to look at almost ten million joint moves.
A good heuristic keeps the search away from most of this space, but each
expansion still pays for the full branching factor, and that is what kills
joint \astar in practice.\index{joint state space!size}\index{branching factor!joint}

\begin{table}[tb]
  \centering\small
  \caption{Size of the joint state space. $|V|^k$ counts all configurations,
  $|V|^{\underline{k}}$ the collision-free ones, and $5^k$ the joint moves of
  one configuration on a 4-connected grid with a wait action (from
  \code{gen\_ch11\_state\_space.py}).}
  \label{tab:ch11-sizes}
  \begin{tabular}{@{}rrrrr@{}}
  \toprule
  $k$ & $54^k$ ($8\times 8$ map) & $54^{\underline{k}}$ & $5^k$ joint moves & $400^k$ ($20\times 20$ map)\\
  \midrule
   1 & $54$ & $54$ & $5$ & $400$\\
   2 & $2.9 \cdot 10^{3}$ & $2.9 \cdot 10^{3}$ & $25$ & $1.6 \cdot 10^{5}$\\
   3 & $1.6 \cdot 10^{5}$ & $1.5 \cdot 10^{5}$ & $125$ & $6.4 \cdot 10^{7}$\\
   4 & $8.5 \cdot 10^{6}$ & $7.6 \cdot 10^{6}$ & $625$ & $2.6 \cdot 10^{10}$\\
   5 & $4.6 \cdot 10^{8}$ & $3.8 \cdot 10^{8}$ & $3\,125$ & $1.0 \cdot 10^{13}$\\
   8 & $7.2 \cdot 10^{13}$ & $4.2 \cdot 10^{13}$ & $3.9 \cdot 10^{5}$ & $6.6 \cdot 10^{20}$\\
  10 & $2.1 \cdot 10^{17}$ & $8.7 \cdot 10^{16}$ & $9.8 \cdot 10^{6}$ & $1.0 \cdot 10^{26}$\\
  20 & $4.5 \cdot 10^{34}$ & $7.8 \cdot 10^{32}$ & $9.5 \cdot 10^{13}$ & $1.1 \cdot 10^{52}$\\
  \bottomrule
  \end{tabular}
\end{table}

\subsection{The heuristic}

Let $d_i(u)$ be the shortest-path distance in $G$ from cell $u$ to the goal
$\gamma_i$ of agent $i$, computed once per agent by a backward breadth-first
search (the backward Dijkstra of \cref{ch:ch03} with unit costs). The
standard heuristic for the joint search is the sum of the individual
distances,\index{heuristic!sum of individual costs}
\begin{equation}
  \hcost(v) = \sum_{i=1}^{k} d_i(v^i).
  \label{eq:ch11-heuristic}
\end{equation}

\begin{proposition}[The sum heuristic is consistent]
\label{thm:ch11-heuristic}
With the costs of \cref{eq:ch11-cost}, the heuristic \cref{eq:ch11-heuristic}
is consistent on the joint graph, hence admissible, and $\hcost(v_f) = 0$.
\end{proposition}
\begin{proof}
$\hcost(v_f) = \sum_i d_i(\gamma_i) = 0$. For a legal joint move from $v$ to
$v'$ and an agent $i$ there are three cases. If $v'^i$ is a neighbour of
$v^i$, then $d_i(v^i) \le 1 + d_i(v'^i)$ by the triangle inequality of
shortest-path distances, and $c_i = 1$. If the agent waits away from its goal,
$d_i(v^i) = d_i(v'^i)$ and $c_i = 1$. If it rests at its goal, $d_i = 0$ on
both sides and $c_i = 0$. In every case $d_i(v^i) \le c_i(v^i, v'^i) +
d_i(v'^i)$; summing over $i$ gives $\hcost(v) \le c(v, v') + \hcost(v')$,
which is \cref{eq:ch04-consistent} of \cref{ch:ch04} on the joint graph.
Admissibility follows from \cref{thm:ch04-consistent-admissible}.
\end{proof}

\subsection{Individual policies and the policy path}

\begin{definition}[Individual policy]
\label{def:ch11-policy}
An \textbf{individual policy}\index{individual policy} for agent $i$ is a
function $\phi_i \colon V \to V$ with $\phi_i(\gamma_i) = \gamma_i$ and, for
every other cell $u$, $\phi_i(u) \in \Succ(u)$ with $d_i(\phi_i(u)) =
d_i(u) - 1$. Following $\phi_i$ from any cell $u$ reaches $\gamma_i$ after
exactly $d_i(u)$ moves and then waits there for free. The \textbf{policy
successor} of a joint configuration $v$ is $\phi(v) = (\phi_1(v^1), \dots,
\phi_k(v^k))$, and the \textbf{policy path}\index{policy path} from $v$ is
$v, \phi(v), \phi(\phi(v)), \dots$ until $v_f$.
\end{definition}

A policy is an optimal single-agent plan in feedback form: it tells the agent
what to do from \emph{every} cell, not only from its start, which is what
\mstar needs when a detour has displaced the agent. When several neighbours
decrease $d_i$, the code takes the first one in the neighbour order east,
west, north, south of the grid; any fixed rule works. Along the policy path
$\fcost = \gcost + \hcost$ stays constant, because every move that costs one
lowers the heuristic by one; the policy path from $v$ costs exactly
$\hcost(v)$ and is therefore an optimal continuation whenever it is
collision-free.

\subsection{Collision sets and limited neighbours}

\begin{definition}[Collision set of a move]
\label{def:ch11-collision}
For a joint move from $v$ to $v'$ in $V^k$ (legal or not), the
\textbf{collision set of the move}\index{collision set} $\Psi(v, v')
\subseteq \set{1, \dots, k}$ is the set of agents that are in a vertex
conflict in $v'$ or in a swap conflict along the move. The move is legal if
and only if $\Psi(v, v') = \emptyset$.
\end{definition}

\begin{definition}[Collision set of a node and limited neighbours]
\label{def:ch11-limited}
During the search every generated configuration $v$ carries a
\textbf{collision set} $C(v) \subseteq \set{1, \dots, k}$, initially empty,
which only ever grows, and a \textbf{back set}\index{back set}
$\mathit{back}(v)$ of all configurations from which $v$ has been generated.
The \textbf{limited neighbours}\index{limited neighbours} of $v$ are the
configurations
\begin{equation}
  N_{C(v)}(v) = \set{v' \in V^k \given v'^i = \phi_i(v^i) \text{ for } i \notin C(v),\
                      v'^i \in \Succ(v^i) \cup \set{v^i} \text{ for } i \in C(v)}.
  \label{eq:ch11-limited}
\end{equation}
\end{definition}

Agents outside the collision set take their policy step; agents inside it
branch over all their moves. With $C(v) = \emptyset$ the only limited
neighbour is the policy successor $\phi(v)$; with $C(v) = \set{1,\dots,k}$
the limited neighbours are all $(b+1)^k$ joint moves. The number of limited
neighbours is $(b+1)^{|C(v)|}$, and $|C(v)|$ is the local dimension of the
search.

% ---------------------------------------------------------------------
\section{The algorithms}
\label{sec:ch11-algorithm}

\subsection{Joint \astar and operator decomposition}
\label{sec:ch11-jointastar}

\textbf{Joint \astar}\index{joint A*} is \cref{alg:ch04-astar} of \cref{ch:ch04}
run on the joint graph $G^k$ with the costs \cref{eq:ch11-cost} and the
heuristic \cref{eq:ch11-heuristic}: the successors of a node are all legal
joint moves, generated as the Cartesian product of the individual move lists
and filtered by \cref{def:ch11-collision}. Since the heuristic is consistent,
\cref{thm:ch04-optimal,thm:ch04-consistent} apply verbatim: joint \astar is
complete, returns an optimal plan, and expands every configuration at most
once. It is the baseline against which everything else in this chapter is
measured, and \cref{tab:ch11-sizes} says why it is only a baseline.

Standley's \textbf{operator decomposition}\index{operator decomposition}
(OD) attacks the branching factor~\cite{standley2010finding}. Instead of
choosing the moves of all $k$ agents in one step, the search chooses them one
agent at a time: from a full configuration it generates the $b+1$ moves of
agent~1, which yield \emph{intermediate} states in which agent~1 has moved and
the others have not; from an intermediate state it generates the moves of the
next agent, checking vertex and swap conflicts only against the agents that
have already chosen; after the move of agent~$k$ a full configuration is
reached again. A node then has at most $b+1$ successors instead of
$(b+1)^k$. The price is a search tree that is $k$ times deeper and a
heuristic that must be evaluated on intermediate states, where it is the sum
of $d_i$ over the new cell of the agents that have moved and the old cell of
the others. Because the intermediate states carry partial $\gcost$ values,
\astar prunes an agent's bad move before the moves of the remaining agents
are ever enumerated; in the experiment of \cref{sec:ch11-example} OD
generates about $200$ times fewer successor candidates than plain joint
\astar with six agents. Its optimality follows from that of \astar, since OD
is \astar on a graph that has the same paths between full configurations.

\subsection{Independence detection}
\label{sec:ch11-id}

Independence detection\index{independence detection} (ID), also due to
Standley~\cite{standley2010finding}, is not a search but a wrapper around
one. \Cref{alg:ch11-id} starts with every agent in a group of its own, plans
each group optimally and independently, and then looks for a pair of groups
whose plans conflict. Such a pair is merged and re-planned as one group with
the joint search; the loop repeats until no two group plans conflict. The
final plan is the union of the group plans.

\begin{algorithm}[htb]
\caption{Simple independence detection}
\label{alg:ch11-id}
\KwIn{graph $G$, starts $s_1, \dots, s_k$, goals $\gamma_1, \dots, \gamma_k$, an optimal joint solver \MstarJointAStar}
\KwOut{an optimal conflict-free plan, or \emph{failure}}
\Fn{\MstarIndependenceDetection{$G$, $s_{1..k}$, $\gamma_{1..k}$}}{
  $\mathcal{G} \gets \set{\set{1}, \set{2}, \dots, \set{k}}$\label{alg:ch11-id:init}\;
  \ForEach{group $A \in \mathcal{G}$}{$\pi_A \gets$ \MstarJointAStar{$G$, $A$}\label{alg:ch11-id:solo}\tcp*{single-agent search}}
  \While{there are groups $A \ne B$ in $\mathcal{G}$ whose plans $\pi_A$ and $\pi_B$ conflict\label{alg:ch11-id:check}}{
    $M \gets A \cup B$; remove $A$ and $B$ from $\mathcal{G}$; add $M$\label{alg:ch11-id:merge}\;
    $\pi_M \gets$ \MstarJointAStar{$G$, $M$}\label{alg:ch11-id:replan}\tcp*{joint search over the merged group only}
    \lIf{$\pi_M = $ failure}{\KwRet \emph{failure}\label{alg:ch11-id:fail}}
  }
  \KwRet $\bigcup_{A \in \mathcal{G}} \pi_A$\label{alg:ch11-id:union}\;
}
\end{algorithm}

\paragraph{Walkthrough.} Line~\ref{alg:ch11-id:solo} runs $k$ single-agent
searches; on a grid they are the \astar searches of \cref{ch:ch04} and cost
almost nothing. Line~\ref{alg:ch11-id:check} is the conflict detection of
\cref{ch:ch07} between the paths of two groups, with the stay-at-target
convention: an agent that has arrived blocks its goal cell for ever. Merging
(line~\ref{alg:ch11-id:merge}) is irreversible, so the groups only grow, and
the loop runs at most $k - 1$ times. Line~\ref{alg:ch11-id:replan} is the
expensive step: a joint search over $|M|$ agents. In the worst case all
agents end up in one group and ID has done $k - 1$ useless searches before the
full joint search; in the common case the largest group stays small
(\cref{fig:ch11-expansions}, right). Standley adds two refinements that keep
groups small: before merging $A$ and $B$, try to re-plan $A$ with the same
cost while treating $\pi_B$ as an obstacle, then the other way round, and
merge only if both attempts fail; and break ties among equally cheap paths in
favour of paths that conflict with fewer other groups, using a
\emph{conflict avoidance table}\index{conflict avoidance table} of the
current plans. The code of this chapter implements the simple version.

\begin{proposition}[Independence detection is optimal]
\label{thm:ch11-id}
If the group solver returns optimal plans and reports failure only when a
group has no plan, then \cref{alg:ch11-id} returns a conflict-free plan of
minimum cost whenever one exists.
\end{proposition}
\begin{proof}
When the loop stops, no two group plans conflict, so their union is a
conflict-free plan; its cost is $\sum_A \cost(\pi_A)$ with each $\pi_A$
optimal for the group $A$. Let $\pi^*$ be any conflict-free plan of all $k$
agents. Restricting $\pi^*$ to the agents of a group $A$ gives a conflict-free
plan for $A$, whose cost is at least $\cost(\pi_A)$ because $\pi_A$ is optimal
for $A$. The cost of $\pi^*$ is the sum of the costs of its restrictions, so
$\cost(\pi^*) \ge \sum_A \cost(\pi_A)$. Hence the returned plan is optimal.
If a merged group has no plan, the whole instance has none, since a plan for
all agents restricts to a plan for the group.
\end{proof}

\subsection{\mstar}
\label{sec:ch11-mstar}

\Cref{alg:ch11-mstar} is the basic \mstar of Wagner and
Choset~\cite{wagner2011mstar,wagner2015subdimensional}. It is joint \astar
with three changes. The successors of a node are its limited neighbours
\cref{eq:ch11-limited} rather than all joint moves. A successor that collides
is not added to the queue; instead its collision set is added to the
collision set of the node and \emph{backpropagated} to every node in the back
set, recursively. A node whose collision set grows is put back into \Open so
that it is expanded again, now with more neighbours.

\begin{algorithm}[htb]
\caption{\mstar with collision sets and backpropagation}
\label{alg:ch11-mstar}
\KwIn{graph $G$, joint start $v_s$, joint goal $v_f$, optimal individual policies $\phi_1, \dots, \phi_k$, heuristic $\hcost$ from \cref{eq:ch11-heuristic}}
\KwOut{an optimal conflict-free plan, or \emph{failure}}
\Fn{\MstarSearch{$G$, $v_s$, $v_f$, $\phi$}}{
  for every configuration (implicitly): $\gcost \gets \infty$, $C \gets \emptyset$, $\mathit{back} \gets \emptyset$\;
  $\gcost(v_s) \gets 0$; \Open $\gets$ priority queue holding $v_s$ with key $(\hcost(v_s), 0)$\label{alg:ch11-mstar:init}\;
  \While{\Open $\ne \emptyset$}{
    $v \gets$ pop the entry with the smallest key $(\fcost, -\gcost)$ from \Open, skipping stale entries\label{alg:ch11-mstar:pop}\;
    \lIf{$v = v_f$}{\KwRet \MstarPath{$v_f$}\label{alg:ch11-mstar:goal}}
    \ForEach{$v' \in$ \MstarNeighbours{$v$, $C(v)$}\label{alg:ch11-mstar:succ}}{
      add $v$ to $\mathit{back}(v')$\label{alg:ch11-mstar:back}\;
      $\Psi \gets$ \MstarCollide{$v$, $v'$}\label{alg:ch11-mstar:collide}\;
      \If{$\Psi \ne \emptyset$}{
        $C(v') \gets C(v') \cup \set{\text{agents sharing a cell in } v'}$\label{alg:ch11-mstar:store}\;
        \MstarBackprop{$v$, $\Psi \cup C(v')$}; \KwContinue\label{alg:ch11-mstar:bpcol}\tcp*{$v'$ is never queued}
      }
      \MstarBackprop{$v$, $C(v')$}\label{alg:ch11-mstar:bpchild}\tcp*{$v'$ may already carry a collision set}
      \If{$\gcost(v) + c(v, v') < \gcost(v')$}{
        $\gcost(v') \gets \gcost(v) + c(v, v')$; parent$(v') \gets v$; push $v'$ into \Open with key $(\gcost(v') + \hcost(v'), -\gcost(v'))$\label{alg:ch11-mstar:relax}\;
      }
    }
  }
  \KwRet \emph{failure}\;
}
\Fn{\MstarNeighbours{$v$, $C$}}{
  \KwRet the Cartesian product of the lists $L_1, \dots, L_k$ with $L_i = \Succ(v^i) \cup \set{v^i}$ if $i \in C$ and $L_i = \set{\phi_i(v^i)}$ otherwise\label{alg:ch11-mstar:limited}\;
}
\Fn{\MstarBackprop{$v$, $C'$}}{
  \If{$C' \not\subseteq C(v)$\label{alg:ch11-mstar:grow}}{
    $C(v) \gets C(v) \cup C'$\;
    \lIf{$v \notin$ \Open \KwAnd $\gcost(v) < \infty$}{push $v$ into \Open again with its current key\label{alg:ch11-mstar:reopen}}
    \lForEach{$u \in \mathit{back}(v)$}{\MstarBackprop{$u$, $C(v)$}\label{alg:ch11-mstar:recurse}}
  }
}
\end{algorithm}

\paragraph{Walkthrough.} Lines~\ref{alg:ch11-mstar:init}
to~\ref{alg:ch11-mstar:goal} are \astar with the key $(\fcost, -\gcost)$ of
\cref{ch:ch04}; the goal test happens when $v_f$ is popped.
Line~\ref{alg:ch11-mstar:succ} is the first difference: with an empty
collision set the loop body runs once, for the policy successor.
Line~\ref{alg:ch11-mstar:back} records \emph{every} parent of $v'$, not only
the cheapest one; backpropagation must reach every node from which a
collision can be reached along limited neighbours, and a node may have
several such parents. Line~\ref{alg:ch11-mstar:collide} computes the
collision set of the move with \cref{def:ch11-collision}. If the move is
illegal, the agents that share a cell in $v'$ are stored in $C(v')$
(line~\ref{alg:ch11-mstar:store}): that set is a property of the
configuration $v'$, which can never be entered, whatever the parent. A swap
conflict, by contrast, is a property of the move and is not stored on $v'$,
because $v'$ may be a perfectly legal configuration when reached by other
moves; it is only backpropagated (line~\ref{alg:ch11-mstar:bpcol}).
Line~\ref{alg:ch11-mstar:bpchild} handles a legal successor that already
carries a collision set from an earlier visit: the parent inherits it, so
that a node always knows about every collision that has been found
downstream of it. Line~\ref{alg:ch11-mstar:relax} is the relaxation of
\astar. In \MstarBackprop, line~\ref{alg:ch11-mstar:grow} stops the
recursion as soon as nothing new arrives, which bounds the total work of
backpropagation by the number of (node, agent) pairs. Line~\ref{alg:ch11-mstar:reopen}
is the re-expansion rule: a node that has already been expanded, but whose
collision set has just grown, goes back into \Open with its old key; when it
is popped again, \MstarNeighbours returns more successors than before, and
the search widens around it. Nodes that are already in \Open need nothing:
they will be expanded with their current, larger collision set when their
turn comes. Line~\ref{alg:ch11-mstar:recurse} pushes the news to all parents.

Two invariants are worth keeping in mind. First, $C(u) \supseteq C(v)$ for
every $u \in \mathit{back}(v)$ at any moment after the propagation has run:
a parent knows at least what its children know. Second, along the policy path
from any expanded node the $\fcost$ value is constant, so the policy path is
explored to its end or to its first collision before any node with a larger
$\fcost$ is touched. \Cref{sec:ch11-properties} builds the correctness proof
on these two facts.

% ---------------------------------------------------------------------
\section{A worked example}
\label{sec:ch11-example}

\begin{example}[Three drones, two aisles, one bay]
\label{ex:ch11-aisles}
\Cref{fig:ch11-example}(a) shows a $5 \times 4$ map with two aisles. The
lower aisle is the row $y = 1$; it has a bay at $(1, 0)$. The upper aisle is
the row $y = 3$; the two are joined only through the passage $(4, 2)$. Drone~1
flies from $(0, 1)$ to $(4, 1)$, drone~2 from $(4, 1)$ to $(0, 1)$, and
drone~3 from $(0, 3)$ to $(4, 3)$. Each drone is four moves from its goal, so
$\hcost(v_s) = 4 + 4 + 4 = 12$. Drones~1 and~2 head straight at each other,
and their policy paths collide in $(2, 1)$ at $t = 2$. Drone~3 never meets
anyone. The instance is \code{worked\_example()} in \code{ch11\_mstar.py};
every number below is produced by its self-test.
\end{example}

\begin{figure}[tb]
  \centering
  \inputfigure{ch11/example}
  \caption{\Cref{ex:ch11-aisles}. (a) The individual policies (arrows with
  arrival times) collide in cell $(2,1)$ at $t = 2$; drone~3 is
  independent. (b) The plan returned by \mstar, of cost $7 + 4 + 4 = 15$:
  drone~1 waits one step at $(1,1)$, ducks into the bay while drone~2 passes
  through $(1,1)$, and continues; drones~2 and~3 follow their policies
  unchanged. The plan is also optimal for the sum of costs and has
  makespan~7.}
  \label{fig:ch11-example}
\end{figure}

\Cref{tab:ch11-trace} lists every expansion of \cref{alg:ch11-mstar} on the
instance, and \cref{fig:ch11-search-graph} draws the search graph. The run
has four phases.

\paragraph{Phase 1: along the policies (steps 1--2).} The start has an
empty collision set and a single limited neighbour, the policy successor
$((1,1), (3,1), (1,3))$ with $\fcost = 12$. That node is expanded next
(step~2), and its policy successor $((2,1), (2,1), (2,3))$ puts drones~1
and~2 in the same cell. The successor is discarded, the set $\set{1,2}$ is
stored on the illegal configuration and backpropagated: first to the node of
step~2, whose collision set becomes $\set{1,2}$, then along its back set to
the start. Both nodes have already been expanded, so both go back into
\Open with their old keys ($\fcost = 12$).

\paragraph{Phase 2: widening (steps 3--4).} The node of step~2 is popped
first because it has the larger $\gcost$. With $C = \set{1,2}$ it now has
$4 \times 3 = 12$ limited neighbours: drone~1 at $(1,1)$ can go east, west,
south into the bay or wait, drone~2 at $(3,1)$ can go west, east or wait,
and drone~3 keeps its policy step to $(2,3)$. One of the twelve is the
collision already known; the other eleven are new nodes with $\fcost = 13$
or $14$. Step~4 re-expands the start with $2 \times 3 = 6$ limited
neighbours, five of them new. Note what has \emph{not} happened: drone~3 is
in no collision set, so it never branches; every node of the search has
drone~3 exactly where its policy puts it. Its dimension has been folded
away.

\paragraph{Phase 3: the patch (steps 5--27).} The nodes with $\fcost = 13$
are those in which drones~1 and~2 have together lost one step against the
heuristic, for instance because one of them waited. Each is expanded with an
empty collision set first, and its policy successor is a swap or a vertex
collision; the collision set grows, the node is re-opened, and it is
expanded again with $C = \set{1,2}$ (steps~5--8). Twelve of the 31
expansions are such second expansions of a node whose collision set grew
after its first expansion. The nodes of steps~9--12 and 21--24 show the other
way a node acquires a collision set: their policy successor is a node that
already carries $\set{1,2}$, and line~\ref{alg:ch11-mstar:bpchild} hands it
up. At $\fcost = 14$ the search tries the bay at the wrong time and lets
drone~2 wander into the passage $(4,2)$ and even to $(4,3)$; none of this
leads anywhere cheaper than 15.

\paragraph{Phase 4: one-dimensional again (steps 28--31).} At $\fcost = 15$
the queue holds the nodes with $\gcost = 9$ generated in steps~7 and~8. In
the node of step~28 drone~1 sits in the bay and drone~2 in $(1,1)$; the
collision set of this node is empty, and its policy successor
$((1,1), (0,1), (4,3))$ is legal: drone~1 steps back into the aisle as
drone~2 leaves it, a following move. From here on every node has an empty
collision set and a single successor; drones~2 and~3 rest at their goals for
free, and drone~1 walks east. The goal is generated in step~31 with
$\gcost = 15$ and popped next, before the three nodes of steps~25--27 that
were re-opened at $\fcost = 15$ with a smaller $\gcost$: the tie-breaking
rule $(\fcost, -\gcost)$ of \cref{ch:ch04} makes the search dive to the goal
instead of widening those nodes first. The returned plan is
\cref{fig:ch11-example}(b).

\begin{table}[tb]
  \centering\footnotesize
  \caption{Every expansion of \cref{alg:ch11-mstar} on \cref{ex:ch11-aisles}.
  A node is written as the cells of drones~1, 2, 3; $C$ is its collision
  set at the moment of the expansion; $\phi(v)$ is its policy successor;
  ``new'' counts successors that enter \Open with a better $\gcost$;
  ``known'' means that the policy successor already had that $\gcost$. The
  $\fcost$ column never decreases.}
  \label{tab:ch11-trace}
  \setlength{\tabcolsep}{3.5pt}
  \begin{tabular}{@{}rlllrrrll@{}}
  \toprule
  Step & Drone 1 & Drone 2 & Drone 3 & $\gcost$ & $\hcost$ & $\fcost$ & $C$ & What happens\\
  \midrule
   1 & $(0,1)$ & $(4,1)$ & $(0,3)$ &  0 & 12 & 12 & $\emptyset$ & policy successor generated\\
   2 & $(1,1)$ & $(3,1)$ & $(1,3)$ &  3 &  9 & 12 & $\emptyset$ & $\phi(v) = (2,1)(2,1)(2,3)$ collides: $C \gets \set{1,2}$,\\
     &         &         &         &    &    &    &             & \quad backpropagated to step~1; both re-opened\\
   3 & $(1,1)$ & $(3,1)$ & $(1,3)$ &  3 &  9 & 12 & $\set{1,2}$ & re-expanded: 11 new, 1 known collision\\
   4 & $(0,1)$ & $(4,1)$ & $(0,3)$ &  0 & 12 & 12 & $\set{1,2}$ & re-expanded: 5 new\\
   5 & $(1,1)$ & $(2,1)$ & $(2,3)$ &  6 &  7 & 13 & $\emptyset$ & $\phi(v)$ swaps 1 and 2: $C \gets \set{1,2}$, re-opened\\
   6 & $(2,1)$ & $(3,1)$ & $(2,3)$ &  6 &  7 & 13 & $\emptyset$ & $\phi(v)$ swaps 1 and 2: $C \gets \set{1,2}$, re-opened\\
   7 & $(1,1)$ & $(2,1)$ & $(2,3)$ &  6 &  7 & 13 & $\set{1,2}$ & re-expanded: 9 new, 3 collisions\\
   8 & $(2,1)$ & $(3,1)$ & $(2,3)$ &  6 &  7 & 13 & $\set{1,2}$ & re-expanded: 3 new, 3 collisions\\
   9 & $(0,1)$ & $(3,1)$ & $(1,3)$ &  3 & 10 & 13 & $\emptyset$ & $\phi(v)$ known; inherits $C = \set{1,2}$ from it\\
  10 & $(1,1)$ & $(4,1)$ & $(1,3)$ &  3 & 10 & 13 & $\emptyset$ & $\phi(v)$ known; inherits $C = \set{1,2}$\\
  11 & $(0,1)$ & $(3,1)$ & $(1,3)$ &  3 & 10 & 13 & $\set{1,2}$ & re-expanded: nothing new\\
  12 & $(1,1)$ & $(4,1)$ & $(1,3)$ &  3 & 10 & 13 & $\set{1,2}$ & re-expanded: 4 new (drone~2 into the passage)\\
  13 & $(1,1)$ & $(3,1)$ & $(2,3)$ &  6 &  8 & 14 & $\emptyset$ & $\phi(v)$ collides in $(2,1)$: re-opened\\
  14 & $(2,1)$ & $(4,1)$ & $(2,3)$ &  6 &  8 & 14 & $\emptyset$ & collides in $(3,1)$: re-opened\\
  15 & $(0,1)$ & $(2,1)$ & $(2,3)$ &  6 &  8 & 14 & $\emptyset$ & collides in $(1,1)$: re-opened\\
  16 & $(1,0)$ & $(2,1)$ & $(2,3)$ &  6 &  8 & 14 & $\emptyset$ & collides in $(1,1)$: re-opened\\
  17 & $(1,1)$ & $(3,1)$ & $(2,3)$ &  6 &  8 & 14 & $\set{1,2}$ & re-expanded: 2 new\\
  18 & $(2,1)$ & $(4,1)$ & $(2,3)$ &  6 &  8 & 14 & $\set{1,2}$ & re-expanded: 3 new\\
  19 & $(0,1)$ & $(2,1)$ & $(2,3)$ &  6 &  8 & 14 & $\set{1,2}$ & re-expanded: nothing new\\
  20 & $(1,0)$ & $(2,1)$ & $(2,3)$ &  6 &  8 & 14 & $\set{1,2}$ & re-expanded: nothing new\\
  21 & $(0,1)$ & $(4,1)$ & $(1,3)$ &  3 & 11 & 14 & $\emptyset$ & $\phi(v)$ known; inherits $C$\\
  22 & $(1,1)$ & $(4,2)$ & $(1,3)$ &  3 & 11 & 14 & $\emptyset$ & $\phi(v)$ known; inherits $C$\\
  23 & $(0,1)$ & $(4,1)$ & $(1,3)$ &  3 & 11 & 14 & $\set{1,2}$ & re-expanded: nothing new\\
  24 & $(1,1)$ & $(4,2)$ & $(1,3)$ &  3 & 11 & 14 & $\set{1,2}$ & re-expanded: 4 new (drone~2 up to $(4,3)$)\\
  25 & $(1,1)$ & $(2,1)$ & $(3,3)$ &  9 &  6 & 15 & $\emptyset$ & $\phi(v)$ swaps 1 and 2: re-opened\\
  26 & $(2,1)$ & $(3,1)$ & $(3,3)$ &  9 &  6 & 15 & $\emptyset$ & $\phi(v)$ swaps 1 and 2: re-opened\\
  27 & $(0,1)$ & $(1,1)$ & $(3,3)$ &  9 &  6 & 15 & $\emptyset$ & $\phi(v)$ swaps 1 and 2: re-opened\\
  28 & $(1,0)$ & $(1,1)$ & $(3,3)$ &  9 &  6 & 15 & $\emptyset$ & drone~1 in the bay: $\phi(v) = (1,1)(0,1)(4,3)$, $\gcost = 12$\\
  29 & $(1,1)$ & $(0,1)$ & $(4,3)$ & 12 &  3 & 15 & $\emptyset$ & drones 2, 3 rest for free: successor has $\gcost = 13$\\
  30 & $(2,1)$ & $(0,1)$ & $(4,3)$ & 13 &  2 & 15 & $\emptyset$ & successor with $\gcost = 14$\\
  31 & $(3,1)$ & $(0,1)$ & $(4,3)$ & 14 &  1 & 15 & $\emptyset$ & generates $v_f$ with $\gcost = 15$; $v_f$ is popped next\\
  \bottomrule
  \end{tabular}
\end{table}

\begin{figure}[tb]
  \centering
  \inputfigure{ch11/search-graph}
  \caption{The search graph of \mstar on \cref{ex:ch11-aisles}, drawn over
  the time axis. Along the policies it is a chain; the collision at $t = 2$
  is backpropagated (red) to the two nodes before it, which are re-expanded
  with the joint moves of drones~1 and~2 (fans of small nodes; only the
  successor on the optimal plan is drawn in full). Once drone~1 has left the
  bay behind drone~2, the collision sets are empty again and the graph is a
  chain to the goal.}
  \label{fig:ch11-search-graph}
\end{figure}

\paragraph{The four searches compared.} On the same instance plain joint
\astar expands 24 nodes but examines $612$ successor candidates, because
every expansion enumerates $5^3 = 125$ joint moves; operator decomposition
expands $100$ states, most of them intermediate, and examines $292$
candidates; independence detection runs three single-agent searches, finds
the conflict between drones~1 and~2, merges them and solves the pair in the
joint space, for $27$ expansions and $162$ candidates in total; \mstar
expands $31$ nodes and examines $127$ candidates. \mstar expands \emph{more}
nodes than joint \astar here: twelve of its expansions are repetitions
forced by collision sets that grew after the first expansion. What it saves
is the branching factor, and on three agents the saving is modest. Of the
$55$ configurations that the search touched, $19$ end with the collision set
$\set{1,2}$ and none contains drone~3.

\subsection{Effort on random maps}
\label{sec:ch11-experiment}

\Cref{fig:ch11-expansions} scales the comparison up. The script
\code{gen\_ch11\_expansions.py} draws $30$ random $8 \times 8$ maps with
$15\,\%$ obstacles for each $k = 2, \dots, 6$, places $k$ agents at random,
and runs the four searches with a cap of $1.5$ million successor
candidates. The left panel shows the median number of candidates. Joint
\astar grows from $140$ for two agents to $77\,000$ for six, roughly a factor
five per agent, as the branching factor predicts; operator decomposition
stays below $360$; independence detection and \mstar stay near $10$ to
$300$ on the median instance, because on the median instance at most two or
three agents interact. The means tell a different story: for $k = 6$ they
are $132\,000$ (joint \astar), $1\,100$ (OD), $14\,000$ (ID) and $42\,000$
(\mstar), dominated by the few instances on which almost every agent ends up
coupled and ID and \mstar degenerate into a joint search plus overhead. The
right panel shows the size of the largest group that ID had to merge and the
largest collision set that \mstar built, on average $2.8$ and $3.8$ agents
out of six. \mstar couples more than ID: it couples agents whenever their
\emph{policy} paths collide anywhere in the region it explores, whereas ID
merges groups only when their optimal plans conflict. Both stay far below
the diagonal $k$, and that gap is the whole point.

\begin{figure}[tb]
  \centering
  \inputfigure{ch11/expansions}
  \caption{Random $8 \times 8$ maps, $30$ instances per $k$. Left: median
  number of successor candidates examined (log scale), with the number of
  collision-free joint configurations $54^{\underline{k}}$ as a reference.
  Right: mean size of the largest group merged by independence detection
  and of the largest collision set built by \mstar; the dashed line is
  ``all $k$ agents''.}
  \label{fig:ch11-expansions}
\end{figure}

% ---------------------------------------------------------------------
\section{Properties: termination, completeness, optimality}
\label{sec:ch11-properties}

Throughout this section $G$ is finite, the policies $\phi_i$ are optimal in
the sense of \cref{def:ch11-policy}, the heuristic is \cref{eq:ch11-heuristic}
and the costs are \cref{eq:ch11-cost}. The results are those of Wagner and
Choset~\cite{wagner2011mstar,wagner2015subdimensional}; the proofs below
follow their argument in the notation of this chapter.

\begin{theorem}[Termination]
\label{thm:ch11-terminates}
\Cref{alg:ch11-mstar} terminates on every instance.
\end{theorem}
\begin{proof}
A node enters \Open for one of two reasons: its $\gcost$ strictly decreased
(line~\ref{alg:ch11-mstar:relax}) or its collision set strictly grew
(line~\ref{alg:ch11-mstar:reopen}). Collision sets are subsets of
$\set{1,\dots,k}$ that only grow, so the second event happens at most $k$
times per node. Between two such events anywhere in the search the graph of
limited neighbours is fixed, and on a fixed finite graph with non-negative
costs the $\gcost$ of a node can decrease only finitely often (the argument
of \cref{thm:ch04-terminates}). There are finitely many configurations,
hence finitely many insertions into \Open, and the loop ends.
\end{proof}

The bound behind this proof is loose: a node can be expanded once per growth
of its collision set and once more per improvement of its $\gcost$. In the
worked example no $\gcost$ ever improved after an expansion, and the
repetitions were exactly the twelve growths.

\begin{theorem}[\mstar is complete and optimal]
\label{thm:ch11-optimal}
If a conflict-free plan exists, \cref{alg:ch11-mstar} returns one of minimum
cost \cref{eq:ch11-cost}. If none exists, it returns \emph{failure}.
\index{M*!completeness}\index{M*!optimality}
\end{theorem}
\begin{proof}[Proof sketch]
Let $C^*$ be the optimal cost. We argue in three steps.

\emph{Step 1: \mstar is \astar on the final limited graph.} Let $G^{\mathrm{sub}}$
be the graph whose vertices are the configurations generated during the run
and whose edges are the limited neighbours \cref{eq:ch11-limited} with
respect to the collision sets at termination. Whenever a collision set grew,
the node was re-inserted, so every node that \mstar expanded was expanded at
least once with its final collision set. \Cref{alg:ch11-mstar} therefore
behaves like \cref{alg:ch04-astar} on $G^{\mathrm{sub}}$, with a consistent
heuristic (\cref{thm:ch11-heuristic}) and with re-opening, and by
\cref{thm:ch04-optimal} it returns a cheapest path of $G^{\mathrm{sub}}$ from
$v_s$ to $v_f$ if $G^{\mathrm{sub}}$ contains one. It remains to show that
$G^{\mathrm{sub}}$ contains a path of cost $C^*$ whenever a plan exists.

\emph{Step 2: policy paths are always explored.} Let $v$ be a node expanded
with $\fcost(v) \le C^*$. Its policy successor $\phi(v)$ is a limited
neighbour whatever $C(v)$ is, and $\fcost(\phi(v)) = \fcost(v)$ by
\cref{def:ch11-policy}. So either $\phi(v)$ is a collision, or it is
generated with $\fcost \le C^*$ and, since the goal is popped only with
$\fcost = $ cost of the returned path and every node with a smaller key is
popped first, it is expanded before the run ends. By induction the policy
path from $v$ is explored up to $v_f$ or up to its first collision, and in
the second case the colliding agents are backpropagated into $C(v)$
(line~\ref{alg:ch11-mstar:bpcol} and the back sets). The same holds for any
path from $v$ along limited neighbours: a collision at its end is
backpropagated to $v$, and $C(u) \supseteq C(w)$ whenever $w$ was generated
from $u$.

\emph{Step 3: no optimal plan is cut off.} Suppose, for contradiction, that
a plan exists but \mstar returns \emph{failure} or a path of cost above
$C^*$. Among all optimal plans $\pi^* = (u_0 = v_s, u_1, \dots, u_m = v_f)$
choose one that maximises the number $j$ of leading nodes $u_0, \dots, u_j$
expanded by \mstar with $\gcost(u_l)$ equal to the cost of the prefix of
$\pi^*$ (every prefix of an optimal path is optimal). Since $v_s$ qualifies,
$j \ge 0$, and $j < m$ because otherwise the goal would have been expanded
with $\gcost = C^*$. Let $A = C(u_j)$ at termination. The move $u_j \to
u_{j+1}$ is not a limited neighbour of $u_j$, otherwise $u_{j+1}$ would have
been generated with optimal $\gcost$ and $\fcost \le C^*$, hence expanded,
contradicting the maximality of $j$. So some agent outside $A$ deviates from
its policy in this move. Now build the path $\sigma$ from $u_j$ in which the
agents in $A$ move exactly as in $\pi^*$ and the agents outside $A$ follow
their policies to their goals. Every agent outside $A$ pays at most what it
pays in $\pi^*$, because a policy path is a shortest path, so
$\cost(\sigma) \le \cost$ of the suffix of $\pi^*$. Two cases.
(i) $\sigma$ is conflict-free. Then the prefix of $\pi^*$ followed by
$\sigma$ is an optimal plan whose first step after $u_j$ \emph{is} a limited
neighbour of $u_j$ (agents outside $A$ follow their policies, agents in $A$
may move freely); that node is generated with optimal $\gcost$, has
$\fcost \le C^*$, and is expanded, so this plan has $j + 1$ leading expanded
nodes, contradicting the choice of $\pi^*$.
(ii) $\sigma$ has a conflict. Its first conflict involves an agent outside
$A$, because the agents in $A$ move as in the conflict-free $\pi^*$. If the
conflict is in the first move, it is detected when $u_j$ is expanded with
$C(u_j) = A$, and backpropagation puts that agent into $A$: a contradiction.
Otherwise the first node $w_1$ of $\sigma$ is a legal limited neighbour of
$u_j$, generated and expanded with $\fcost \le C^*$ (Step~2), and by
backpropagation $C(w_1) \subseteq A$. If $C(w_1) = A$, the next move of
$\sigma$ is a limited neighbour of $w_1$ and the argument repeats one node
further along $\sigma$, with the conflict one step closer. If
$C(w_1) \subsetneq A$, repeat the construction from $w_1$ with the smaller
set $C(w_1)$ in place of $A$: agents in $C(w_1)$ keep following $\sigma$,
the others switch to their policies, and again either the new path is
conflict-free, which gives an optimal plan with a longer expanded prefix, or
its first conflict involves an agent outside $C(w_1)$. Collision sets can
shrink at most $k$ times along this process and, while they stay the same,
the first conflict comes one step closer at every repetition, so after
finitely many repetitions a conflict involving an agent outside the current
collision set is found at a generated successor of an expanded node.
Backpropagation then adds that agent to the collision set of that node and
of all its ancestors on the constructed path, including the node from which
the construction started, contradicting the assumption that the set was
final. Hence $G^{\mathrm{sub}}$ contains an optimal plan, and by Step~1
\mstar returns one. If no plan exists, no path to $v_f$ exists in
$G^{\mathrm{sub}} \subseteq G^k$ either, and by \cref{thm:ch11-terminates}
the algorithm stops with \emph{failure}.
\end{proof}

The proof uses the optimality of the individual policies in one place only,
in case~(i): a policy path is never more expensive than the path it replaces.
With suboptimal policies \mstar stays complete but loses optimality. It uses
the heuristic only through Step~1, so an inflated heuristic $w\hcost$ gives
a $w$-bounded plan by \cref{thm:ch04-weighted}; that is \emph{inflated}
\mstar (\cref{sec:ch11-variants}).

\begin{corollary}[No collisions, no search]
\label{thm:ch11-linear}
If the policy path from $v_s$ is conflict-free, \mstar expands exactly its
nodes, one successor each, and returns it: the number of expansions is the
makespan of the individual paths plus one, whatever $k$ is.
\end{corollary}
\begin{proof}
Every node on the path has an empty collision set and the policy successor as
its only limited neighbour; all have $\fcost = \hcost(v_s)$, no other node
is ever generated, and the goal is reached with $\gcost = \hcost(v_s)$,
which is optimal by admissibility.
\end{proof}

\paragraph{Complexity.} Between the two extremes the effort is that of a
joint search over the agents in the collision sets: a node with $|C(v)| = m$
has $(b+1)^m$ limited neighbours, and the number of such nodes grows with
the size of the region in which those $m$ agents interact. In the worst
case every collision set becomes $\set{1, \dots, k}$ and \mstar is joint
\astar with the extra cost of re-expansions and backpropagation; in the
random-map experiment this happened on the instances that dominate the means
in \cref{sec:ch11-experiment}. The exponent is the size of the largest
collision set, not $k$, which is why Wagner and Choset describe the method as
searching a low-dimensional subspace embedded in the joint space.
\index{M*!complexity}

% ---------------------------------------------------------------------
\section{Constructive methods: Push-and-Swap and Push-and-Rotate}
\label{sec:ch11-constructive}

\subsection{Pebbles instead of paths}

The constructive methods see the agents as \textbf{pebbles}\index{pebble motion}
on the vertices of $G$: one pebble moves at a time, along an edge, into an
empty vertex. This is the \emph{pebble motion} problem studied by Kornhauser,
Miller and Spirakis~\cite{kornhauser1984pebble}, who showed that whether a
given arrangement can be turned into another is decidable in polynomial
time and that, when it can, a sequence of $\bigO{|V|^3}$ moves suffices.
For the special case of a single empty vertex on a graph without a cut
vertex, Wilson~\cite{wilson1974graph} had already characterised the
reachable arrangements: the generalised 15-puzzle. A sequential plan
converts into a MAPF plan by letting every move take one time step during
which all other agents wait, and it is conflict-free by construction. Finding
a \emph{shortest} plan remains NP-hard (\cref{ch:ch07};
\cite{surynek2010optimization,yu2013structure}), so a polynomial method must
give up optimality. Push-and-Swap and Push-and-Rotate do exactly that.

\subsection{The push primitive}

\PasPush{$a$} moves agent $a$ along its shortest path $p$ towards its goal
(\cref{fig:ch11-push-swap}(a)). When the next vertex on $p$ is empty, $a$
steps into it. When it is occupied by an agent $b$ that is not yet
\emph{locked} at its own goal, the algorithm \emph{clears} the vertex: a
breadth-first search from the vertex, avoiding the remaining path $p$ and
the vertices of locked agents, finds the nearest empty vertex $e$, and every
agent on the chain from $b$ to $e$ shifts one step towards $e$, starting
with the agent next to $e$. Then $a$ steps forward. In the figure $b$ and
$c$ shift one step down the side branch, after which $a$ walks through:
$2 + 4 = 6$ single moves in total. The push fails when the blocking agent is
locked, or when no empty vertex can be reached without crossing $p$; in both
cases a swap is needed.\index{push (primitive)}

\begin{figure}[tb]
  \centering
  \inputfigure{ch11/push-swap}
  \caption{The primitives. (a) A push with a chain shift: $c$ moves first,
  then $b$, then $a$ walks through. (b) The exchange at the heart of a swap:
  with $a$ at a vertex $v$ of degree three, $b$ behind it and two empty
  neighbours $n_1, n_2$, six moves exchange the two agents and leave the
  rest of the graph untouched. (c) A rotation: with one cleared vertex next
  to a fully occupied cycle, every agent on the cycle advances one position
  in six moves.}
  \label{fig:ch11-push-swap}
\end{figure}

\subsection{The swap primitive}

\PasSwap{$a$, $b$} exchanges the positions of two adjacent agents without
changing, in the end, the position of anyone else. The exchange itself needs
a vertex $v$ with at least three neighbours and two empty neighbours of $v$:
\cref{fig:ch11-push-swap}(b) shows the six moves. To get there the algorithm
searches for the nearest such vertex in order of distance, and for each
candidate it performs a \emph{multipush}\index{swap (primitive)}: the agent
nearer to $v$ leads along a shortest path to $v$ and the other follows one
step behind, pushing other agents off the path with the clearing rule of the
push. At $v$ two further neighbours are cleared in the same way. Then the six
exchange moves are executed, and finally the multipush and the clearing are
undone in reverse order, move by move, whoever now stands on the vertex
concerned moving back. Because the two agents have exchanged roles, the
reversal brings $b$ to where $a$ started and $a$ to where $b$ started, and
every other agent returns to its old vertex. If the leader is $d$ steps from
$v$ and the vertices on the way and next to $v$ are empty, the swap costs
$2d$ moves out, $6$ moves for the exchange and $2d$ moves back: $4d + 6$
moves, verified for $d = 1, \dots, 4$ by the self-test of the chapter code.
If no candidate vertex allows the manoeuvre, the swap fails and so does the
algorithm.

\subsection{Push-and-Swap}

\Cref{alg:ch11-pas} gives the algorithm of Luna and
Bekris~\cite{luna2011push} at the level of its primitives. Agents are
solved one after another. The current agent pushes until it reaches its goal
or is blocked by a locked agent, swaps with the blocker, and pushes again;
when it arrives it is locked. A locked agent that a swap has displaced is
unlocked and solved again later, which the original paper handles in a
\emph{resolve} step. Since every swap moves the current agent one vertex
closer to its goal along its path and leaves all other agents where they
were, the current agent arrives after finitely many primitives, and the loop
over the agents terminates.

\begin{algorithm}[htb]
\caption{Push-and-Swap (conceptual)}
\label{alg:ch11-pas}
\KwIn{graph $G$, starts $s_1, \dots, s_k$, goals $\gamma_1, \dots, \gamma_k$, with at least two empty vertices}
\KwOut{a sequence of single-agent moves, or \emph{failure}}
\Fn{\PasPlan{$G$, $s_{1..k}$, $\gamma_{1..k}$}}{
  $U \gets \emptyset$\tcp*{locked agents, already at their goals}
  \ForEach{agent $a$ in a fixed order\label{alg:ch11-pas:agent}}{
    \While{$a$ is not at $\gamma_a$}{
      $p \gets$ shortest path from the position of $a$ to $\gamma_a$\;
      \PasPush{$a$, $p$, $U$}\label{alg:ch11-pas:push}\tcp*{as far as possible; clears unlocked blockers}
      \If{$a$ is not at $\gamma_a$\label{alg:ch11-pas:blocked}}{
        $b \gets$ the locked agent on the next vertex of $p$\;
        \lIf{\PasSwap{$a$, $b$} fails}{\KwRet \emph{failure}\label{alg:ch11-pas:swap}}
        remove $b$ from $U$ and schedule it to be solved again\label{alg:ch11-pas:resolve}\;
      }
    }
    add $a$ to $U$\label{alg:ch11-pas:lock}\;
  }
  \KwRet the recorded moves\;
}
\end{algorithm}

Luna and Bekris state that the algorithm is complete for every instance
whose graph has at least two unoccupied vertices, that is, at most
$|V| - 2$ agents on a connected graph: it returns a plan whenever one
exists.\index{Push-and-Swap!completeness} The two empty vertices are what
the exchange of \cref{fig:ch11-push-swap}(b) needs, and a degree-three
vertex exists in every connected graph that is not a path or a cycle. De
Wilde, ter Mors and Witteveen~\cite{dewilde2014push} later found instances
in this class on which the published algorithm fails, in particular when
solved agents must be moved in a coordinated way or when the agents that
need to be exchanged sit on a fully occupied cycle, and they repaired it.

\subsection{Push-and-Rotate}

Push-and-Rotate~\cite{dewilde2014push} keeps push and swap and adds two
ingredients. The first is the \PasRotate primitive for a cycle whose
vertices are all occupied (\cref{fig:ch11-push-swap}(c)): a vertex next to
the cycle is cleared, one agent steps out of the cycle into it, the
remaining agents on the cycle advance one position each into the vertex just
vacated, and the agent that stepped out steps back in behind them. Every
agent on the cycle has now moved one vertex along it, which a push could
never achieve because no vertex of the cycle was empty.\index{rotate (primitive)}
The second ingredient is a decomposition of the graph into
\emph{subproblems}: maximal subgraphs within which the agents assigned to
them can reach every position, found from the biconnected components and
the trees hanging off them. Agents are assigned to subproblems, the
subproblems are ordered by a priority that says which of them must be solved
first, and an instance is recognised as unsolvable when an agent cannot be
brought into the subproblem of its goal. This is where the theory of
Kornhauser, Miller and Spirakis enters the algorithm.

\begin{theorem}[Completeness of Push-and-Rotate; de Wilde, ter Mors and Witteveen 2014]
\label{thm:ch11-par}
For every MAPF instance whose graph has at least two unoccupied vertices,
Push-and-Rotate returns a conflict-free plan if the instance is solvable and
reports failure otherwise.\index{Push-and-Rotate!completeness}
\end{theorem}
\begin{proof}[Proof sketch]
The two empty vertices guarantee that within a biconnected subproblem any
two agents can be exchanged by a swap or a rotation, so the reachable
arrangements of a subproblem are all permutations of its agents; the
decomposition and the priorities reduce the instance to a sequence of such
subproblems in which pushes move agents between them. The full proof, some
thirty pages of case analysis, is in \cite{dewilde2014push}.
\end{proof}

\subsection{Feasibility is not optimality}
\label{sec:ch11-quality}

Both methods are \textbf{feasibility solvers}\index{feasibility solver}: they
never compare the cost of two plans, and nothing in them bounds the cost of
the plan they return. The swap example of \cref{fig:ch11-push-swap}(b) makes
this concrete. On the corridor $c_0 c_1 c_2 c_3 c_4$ with a spur $s$ at $c_3$
(\code{swap\_example()} in the chapter code), agent $a$ starts at $c_0$ with
goal $c_1$ and agent $b$ at $c_1$ with goal $c_0$. The driver of the chapter
code solves it with $14$ single-agent moves: $b$ is pushed aside to $c_2$,
$a$ steps to its goal and is locked, $b$ is blocked by the locked $a$, the
two swap at $c_3$ (a multipush of two moves, six exchange moves, two moves
back: $4 \cdot 1 + 6 = 10$), $b$ walks to $c_0$, and the displaced $a$ is
solved again. Executed one move per
time step this plan has makespan $14$ and sum of costs $27$. The optimal
joint plan, found by joint \astar, has makespan $7$ and sum of costs $14$:
both agents walk to the junction together, $a$ ducks into the spur, and both
walk back. The ratio is two on this toy graph; the term $4d + 6$ of every
swap grows with the distance $d$ to the nearest junction, and a swap is
repeated for every locked agent in the way. Both papers therefore
post-process their plans: moves of different agents that do not interfere
are executed in the same time step, and redundant moves are removed. Even
so, the plans are typically several times longer than optimal on crowded
instances, and no bound is available. What the methods buy with this is
speed on instances that are hopeless for optimal solvers: hundreds of agents
on a graph with only a few free vertices, where \cbs would face an
astronomical constraint tree and \mstar a collision set containing everyone.

% ---------------------------------------------------------------------
\section{Variants and extensions}
\label{sec:ch11-variants}

\paragraph{Recursive \mstar.} Basic \mstar has one weakness that
\cref{fig:ch11-expansions} shows: its collision sets are flat. If drones~1
and~2 collide in one corner of the map and drones~3 and~4 in another, a node
whose descendants see both collisions gets $C = \set{1,2,3,4}$ and $5^4$
limited neighbours, although the two pairs never interact.
\textbf{Recursive \mstar}\index{M*!recursive (rM*)} (r\mstar) keeps the
collision set as a collection of disjoint subsets, $\set{\set{1,2},
\set{3,4}}$, and computes the ``policy'' of each subset by a recursive
\mstar call over the agents of that subset only; the top level then branches
over the product of the subsets' policies instead of over all moves of all
four agents. The cost becomes exponential in the size of the largest subset
that must be coupled, not in the size of their union.
\textbf{ODr\mstar}~\cite{ferner2013odrmstar} adds operator decomposition
inside the sub-searches and was, at the time of its publication, among the
strongest optimal MAPF solvers on the standard benchmarks;
\textcite{wagner2015subdimensional} give the full treatment of r\mstar and
its proofs.

\paragraph{Inflated \mstar.} Replacing $\hcost$ by $w\hcost$ with $w > 1$
in \cref{alg:ch11-mstar} gives \textbf{inflated
\mstar}\index{M*!inflated}: Step~1 of the proof of \cref{thm:ch11-optimal}
now invokes \cref{thm:ch04-weighted} instead of \cref{thm:ch04-optimal}, so
the plan costs at most $w C^*$, and the inflated search reaches the goal
with far fewer expansions inside the coupled patches. It is the
bounded-suboptimal relative of \mstar in the same sense in which \ecbs is
the bounded-suboptimal relative of \cbs.

\paragraph{Relatives among the optimal solvers.} Meta-agent
\cbs~\cite{sharon2015cbs} merges two agents into a \emph{meta-agent},
planned by a joint search, once they have conflicted more than a threshold
number of times in the constraint tree; it is the bridge between \cbs and
independence detection, and ODr\mstar has been used as its meta-agent
solver. The increasing cost tree search and enhanced partial expansion
\astar attack the joint space in yet other ways;
\textcite{felner2017search} survey all of them and the trade-offs between
them. \textcite{standley2011complete} turn independence detection and
operator decomposition into a complete anytime algorithm that returns a
first plan quickly and improves it while time remains.

\paragraph{Relatives among the constructive methods.} Surynek's
BIBOX~\cite{surynek2009novel} solves instances on biconnected graphs with
two empty vertices by a different set of macros, and the polynomial
algorithm of Kornhauser, Miller and Spirakis~\cite{kornhauser1984pebble}
is the theoretical ancestor of all of them. None of these bounds the cost
of the plan; when a bound matters, a bounded-suboptimal search such as
\ecbs or inflated \mstar is the right tool.

\subsection{Every MAPF method of the book in one table}
\label{sec:ch11-comparison}

\Cref{tab:ch11-comparison} places the methods of \cref{ch:ch08,ch:ch09,ch:ch10}
and of this chapter side by side. Read it column by column. Completeness
is the promise to find a plan whenever one exists; only prioritized planning
breaks it, and only Push-and-Rotate also promises to \emph{report} that no
plan exists. Optimality is expensive everywhere: the three optimal methods
pay in the size of their tree or their collision sets. A bound is the middle
ground, and two methods offer it. The ``scales with'' column names the
quantity that actually decides the running time, which is never simply $k$.

\begin{table}[tb]
  \centering\footnotesize
  \caption{The multi-agent path finding methods of this book. ``Bounded''
  means bounded suboptimality with a user-chosen factor $w$; ``scales
  with'' names the quantity that dominates the running time.}
  \label{tab:ch11-comparison}
  \setlength{\tabcolsep}{4pt}
  \begin{tabularx}{\textwidth}{@{}L{2.5cm}YYYYYl@{}}
  \toprule
  Method & Complete? & Optimal? & Bounded? & Scales with & Typical use & Ch.\\
  \midrule
  Prioritized planning & no in general; yes on well-formed instances & no & no & $k$ single-agent space-time searches & very large swarms, real-time fallback & \ref{ch:ch08}\\
  Joint \astar, with OD & yes & yes & -- & $|V|^k$ states; $(b{+}1)^k$ moves per node, $b{+}1$ with OD & baseline, $k \le 4$ & \ref{ch:ch11}\\
  Independence detection & yes, with a complete group solver & yes, with an optimal group solver & inherited & size of the largest conflicting group & wrapper around any solver & \ref{ch:ch11}\\
  \cbs (ICBS) & yes on solvable instances & yes & -- & conflicts resolved: size of the constraint tree & global planner of \cref{ch:ch24}, tens of drones & \ref{ch:ch09}\\
  \ecbs & yes on solvable instances & no & yes, $w$ & conflicts, fewer as $w$ grows & same, hundreds of drones & \ref{ch:ch10}\\
  \mstar (r\mstar, ODr\mstar) & yes & yes; inflated: no & inflated: $w$ & size of the largest collision set & sparse interactions, optimal plans & \ref{ch:ch11}\\
  Push-and-Swap, Push-and-Rotate & yes with $\ge 2$ empty vertices (P\&R); reports unsolvable & no & no & polynomial in $|V|$ and $k$ & dense, puzzle-like instances; any feasible plan & \ref{ch:ch11}\\
  \bottomrule
  \end{tabularx}
\end{table}

% ---------------------------------------------------------------------
\section{Implementation notes}
\label{sec:ch11-implementation}

The file \code{ch11\_mstar.py} contains a small graph class built from an
ASCII map, the individual policies (\code{Policies}), joint \astar with an
optional operator decomposition, simple independence detection, \mstar, a
plan validator and the push-and-swap primitives, all in about 900 lines.
The design decisions that matter are the following.

\paragraph{Nodes and the queue.} Every generated configuration owns a
record with its $\gcost$, parent, collision set (a Python \code{set} of
agent indices; a bit mask is faster for $k > 16$), back set and an
\code{in\_open} flag. The heap holds tuples $(\fcost, -\gcost, \text{counter},
\text{configuration})$ as in \cref{ch:ch04}. A node is pushed when its
$\gcost$ improves and when its collision set grows while it is not in the
queue; an entry popped with a stale $\gcost$, or for a node whose flag says
that it has been expanded since the push, is skipped. The flag is what
makes ``$v \notin$ \Open'' in \cref{alg:ch11-mstar} cheap.

\paragraph{Policies and limited neighbours.} One breadth-first search per
goal gives the distance tables; the next policy step is computed on demand
and cached. \Cref{lst:ch11-neighbours} shows how the limited neighbours are
built as a Cartesian product of per-agent move lists, the wait first so that
ties are broken deterministically, and how the collision sets are
backpropagated with an explicit stack rather than by recursion, which would
overflow on long chains.

\begin{lstlisting}[caption={Limited neighbours and backpropagation (from \code{code/ch11\_mstar.py}; \code{push} and \code{backprop} are local functions of \code{mstar}).},label={lst:ch11-neighbours}]
def limited_neighbours(node, inst, pol):
    """Agents in the collision set may move anywhere; the rest follow their policy."""
    options = []
    for i, v in enumerate(node.config):
        if i in node.collision_set:
            options.append(agent_moves(inst.graph, v))
        else:
            options.append([pol.step(i, v)])
    return itertools.product(*options)

    def push(node):
        node.in_open = True
        heapq.heappush(open_heap, (node.g + pol.h(node.config), -node.g,
                                   next(counter), node.config))

    def backprop(config, colset):
        """Add colset to the collision set of config and of all its ancestors."""
        stack = [(config, colset)]
        while stack:
            cfg, cs = stack.pop()
            node = nodes[cfg]
            if cs <= node.collision_set:
                continue
            node.collision_set |= cs
            if not node.in_open and node.g < INF:
                push(node)                    # re-expand with more neighbours
            for p in node.back_set:
                stack.append((p, node.collision_set))
\end{lstlisting}

\paragraph{The successor loop.} \Cref{lst:ch11-loop} is the body of the
main loop after a node \code{cfg} has been popped. The collision check
\code{collisions} finds vertex collisions with a dictionary from cells to
agents and swap collisions by looking up which agent arrives at the cell an
agent is leaving. A colliding successor is never queued; only the agents
that share a cell in it are stored on it, while the whole collision set of
the move is backpropagated from the parent. The two \code{rec} lines record
the trace that produced \cref{tab:ch11-trace}.

\begin{lstlisting}[caption={The successor loop of \mstar (from \code{code/ch11\_mstar.py}).},label={lst:ch11-loop}]
        for nxt in limited_neighbours(node, inst, pol):
            generated += 1
            child = nodes.get(nxt)
            if child is None:
                child = nodes[nxt] = MstarNode(nxt)
            child.back_set.add(cfg)
            col = collisions(cfg, nxt)
            if col:
                # a vertex collision belongs to the configuration nxt, a swap
                # collision only to the edge cfg -> nxt
                child.collision_set |= vertex_collisions(nxt)
                rec["collisions"].append((nxt, frozenset(col)))
                backprop(cfg, col | child.collision_set)
                continue
            backprop(cfg, child.collision_set)
            new_g = node.g + step_cost(cfg, nxt, goals)
            if new_g < child.g:
                child.g = new_g
                child.parent = cfg
                push(child)
                rec["successors"].append((nxt, new_g))
\end{lstlisting}

\begin{pitfall}[A swap is a property of the move, not of the configuration]
The first version of the chapter code stored the whole collision set of an
illegal move on the successor configuration. For a vertex collision that is
right: the configuration itself is impossible. For a swap it is wrong: the
configuration after the swap is a legal state that the search may reach
later by other moves, and it then dragged a stale collision set into every
node that generated it. On \cref{ex:ch11-aisles} this coupled drones~1
and~2 in nodes where they had long passed each other. The search stayed
correct, because larger collision sets never remove an optimal path, but
it did more work than necessary, and on larger instances the extra coupling
grows. Store vertex collisions on the configuration and swap collisions on
the parent only.
\end{pitfall}

\begin{pitfall}[All parents, and always re-open]
The back set of a node must contain \emph{every} node from which it was
generated, not only the parent of its cheapest path, and a node whose
collision set grows must be pushed back into \Open even though it has
already been expanded. Dropping either shortcut makes \mstar incomplete:
a collision found downstream never reaches the node that needs to branch,
its limited neighbours never include the detour, and the search reports
failure on solvable instances. The self-test guards against this by
comparing the cost of \mstar with joint \astar on random instances.
\end{pitfall}

\begin{pitfall}[Free waits at the goal are not the sum of costs]
The cost convention \cref{eq:ch11-cost} lets an agent rest at its goal for
free, which keeps the cost of a move a function of the two configurations.
The sum of costs of \cref{ch:ch07} charges every step until the agent's
\emph{final} arrival, so an agent that reaches its goal, waits, and later
steps aside is charged for the wait as well. The two objectives agree on
every plan in which nobody leaves a goal, and the chapter code checks that
they agree on all its examples, but a plan that is optimal for one need not
be optimal for the other. Standley's original implementation charges the
free waits retroactively when the agent moves again, which makes the state
depend on history; \cbs sidesteps the problem because its low-level search
plans one agent at a time in space-time and knows the arrival time.
\end{pitfall}

\paragraph{The push-and-swap primitives.} The class \code{PushSwapState}
holds the positions and a log of single-agent moves; \code{clear} shifts a
chain of agents towards the nearest empty vertex found by a breadth-first
search that avoids forbidden vertices, \code{push} walks along a path and
calls \code{clear}, and \code{swap} tries the vertices of degree at least
three in order of distance, performs the multipush, the clearing and the
six exchange moves, and then undoes the preparation \emph{by position}:
for every logged preparatory move, whoever now stands on its target vertex
moves back to its source. This one rule implements the reversal with the
agents exchanged. The driver \code{push\_and\_swap} solves the examples of
the chapter but is not the complete algorithm: it has no resolve step and no
graph decomposition.

% ---------------------------------------------------------------------
\section{Where this fits in the drone system}
\label{sec:ch11-drone}

\begin{dronebox}
\index{M*!in the drone system}%
The global planner of the hybrid architecture in \cref{ch:ch24} is \cbs or
\ecbs, and this chapter explains why. A swarm planner should \emph{couple}
drones only where they interact, and it should do so with a guarantee. \cbs
couples through constraints and never enumerates a joint space; \ecbs adds a
bound that the operator can set. Independence detection is the cheapest way
to make any of these solvers scale: split the swarm into groups whose
nominal plans do not conflict, plan the groups separately, and merge only
on conflict; the largest group, not the swarm, decides the running time.
\mstar is the right global planner when interactions are rare but must be
resolved optimally, for instance for a few dozen drones on a large map with
occasional crossings, and inflated \mstar when a bound is enough; on the
instances of \cref{fig:ch11-expansions} it and independence detection keep
the effort near that of the single-agent searches. A feasible plan is worth
accepting when the alternative is no plan at all: a dense hangar in which
drones must be reshuffled before take-off, or a replanning request that
must be answered within one control cycle, is where Push-and-Rotate's
guarantee earns its place as a fallback. What none of these methods does is
react: they plan for the swarm's own members on a known map, and the
intruder that appears after take-off is the business of the prediction and
local avoidance layers of \cref{ch:ch18,ch:ch12,ch:ch13}. In the study plan
(\cref{ch:appA}) this chapter belongs to Week~5, as the reading that
completes the practical MAPF picture next to \ecbs.
\end{dronebox}

% ---------------------------------------------------------------------
\section{Summary}
\begin{summary}
\begin{itemize}
  \item The joint state space of $k$ agents has $|V|^k$ configurations and
        $(b+1)^k$ joint moves per node; joint \astar with the sum-of-distances
        heuristic is complete and optimal but pays the full branching factor
        at every expansion. Operator decomposition cuts the branching factor
        to $b+1$ by moving one agent at a time through intermediate states.
  \item Independence detection solves groups separately and merges two
        groups only when their optimal plans conflict; with an optimal group
        solver the result is optimal (\cref{thm:ch11-id}), and the running
        time is governed by the largest group.
  \item \mstar follows the individual optimal policies through the joint
        space and branches only over the agents in a node's collision set.
        Collisions found downstream are backpropagated along the back sets,
        and a node whose collision set grew is expanded again. Vertex
        collisions belong to the configuration, swap collisions to the move.
  \item With optimal policies and the consistent sum heuristic, \mstar
        terminates and is complete and optimal (\cref{thm:ch11-optimal});
        without collisions it expands only the policy path
        (\cref{thm:ch11-linear}); in the worst case it is joint \astar plus
        overhead. Recursive, inflated and OD variants exist.
  \item Push-and-Swap and Push-and-Rotate move one agent at a time with the
        push, swap and rotate primitives. Push-and-Rotate is complete for
        instances with at least two empty vertices; both are feasibility
        solvers whose plans can be several times longer than optimal.
  \item \Cref{tab:ch11-comparison} summarises every MAPF method of the book;
        the swarm planner of \cref{ch:ch24} uses \cbs/\ecbs, with
        independence detection and \mstar as exact alternatives when
        interactions are sparse and Push-and-Rotate as a fallback when any
        feasible plan will do.
\end{itemize}
\end{summary}

\paragraph{Further reading.}
The joint-space formulation, operator decomposition and independence
detection are from \textcite{standley2010finding}; \textcite{standley2011complete}
make them complete and anytime. \mstar was introduced by
\textcite{wagner2011mstar} and developed, with recursive and inflated
variants and the full proofs, in the journal version
\cite{wagner2015subdimensional}; ODr\mstar is in \cite{ferner2013odrmstar}.
Push-and-Swap is \cite{luna2011push}; Push-and-Rotate, its completeness
proof and the correction of the earlier claims are in \cite{dewilde2014push};
the pebble-motion theory behind them is \cite{kornhauser1984pebble} and
\cite{wilson1974graph}, and BIBOX is \cite{surynek2009novel}. The
definitions of conflicts and objectives follow \textcite{stern2019mapf}, and
\textcite{felner2017search} survey the search-based optimal solvers,
including \cbs, ICTS, EPEA\textsuperscript{*} and \mstar, with a comparison
of their strengths. \textcite{lavalle2006planning} treats multi-robot
planning in configuration space at textbook level, including the joint
state space of this chapter.

% ---------------------------------------------------------------------
\section{Exercises}
\label{sec:ch11-exercises}

\begin{exercise}[\difficulty{1}]
\label{exr:ch11-sizes}
A warehouse map has $|V| = 300$ free cells and is 4-connected. (a) How many
joint configurations, how many collision-free configurations and how many
joint moves per node does the joint graph have for $k = 3$ and for $k = 12$
agents? (b) Operator decomposition has $b + 1 = 5$ successors per state but
needs $k$ intermediate levels to move all agents once. Compare the number of
states generated by one full joint step of plain joint \astar and of OD for
$k = 12$, assuming no pruning. (c) Explain in two sentences why a good
heuristic does not rescue plain joint \astar from the branching factor.
\end{exercise}

\begin{exercise}[\difficulty{1}]
\label{exr:ch11-collision-sets}
Take the four-vertex graph of \cref{fig:ch11-joint-graph} and add a third
agent that starts at $d$ with goal $d$. Agents~1 and~2 are as in the figure.
Run \cref{alg:ch11-mstar} by hand for the first three expansions: give the
policy of every agent, the nodes expanded, the collisions found, the
collision sets after backpropagation and the limited neighbours of the
re-expanded start. Which agent never enters a collision set, and why not,
although the two-agent plan of cost $7$ in the figure needs the vertex $d$?
What does \mstar return on this instance, and which theorem guarantees that
it returns at all?
\end{exercise}

\begin{exercise}[\difficulty{2}]
\label{exr:ch11-heuristic}
(a) Prove that $\hcost_{\max}(v) = \max_i d_i(v^i)$ is admissible for the
cost \cref{eq:ch11-cost} and that it is dominated by \cref{eq:ch11-heuristic}
in the sense of \cref{def:ch04-dominance}. (b) Suppose waiting were free
everywhere, not only at the goal. Is \cref{eq:ch11-heuristic} still
consistent? Is the joint search still finite? (c) The makespan objective of
\cref{ch:ch07} charges one per time step regardless of how many agents
move. Give a joint move cost and an admissible heuristic for it, and explain
why the policy path is still an optimal continuation when it is
conflict-free.
\end{exercise}

\begin{exercise}[\difficulty{2}]
\label{exr:ch11-id-groups}
(a) Construct an instance with three agents in which simple independence
detection (\cref{alg:ch11-id}) merges all three agents although a plan
exists in which the optimal path of agent~3 conflicts with nobody; use the
fact that a group may have several optimal plans. (b) Explain how Standley's
conflict avoidance table changes the outcome on your instance. (c) Show that
\cref{thm:ch11-id} still holds when the group solver is \mstar or \cbs, and
what happens to the guarantee when it is \ecbs.
\end{exercise}

\begin{exercise}[\difficulty{2}]
\label{exr:ch11-reexpansion}
In \cref{tab:ch11-trace}, twelve nodes are expanded twice. (a) For each of
them, say whether its collision set grew because of a collision among its
own successors or because it was handed up from a child
(line~\ref{alg:ch11-mstar:bpchild}). (b) Explain why the second expansion of
the start node (step~4) could not have been avoided by expanding the start
with $C = \set{1,2}$ right away, that is, why \mstar cannot know the
collision set of a node before it has searched below it. (c) Design an
instance with three drones in one aisle in which the collision set of the
start grows twice, first to $\set{1,2}$ and later to $\set{1,2,3}$, and
verify your prediction with \code{mstar(inst, trace=True)}.
\end{exercise}

\begin{exercise}[\difficulty{2}]
\label{exr:ch11-swap-count}
(a) Prove that the swap of \cref{sec:ch11-constructive} needs exactly
$4d + 6$ single-agent moves when the leading agent is $d$ steps from the
junction and no other agent has to be cleared. (b) Executed one move per
time step, what are the makespan and the sum of costs of the plan for the
swap example of \cref{sec:ch11-quality} as functions of $d$, and what are
those of the optimal joint plan? (c) Show that the six exchange moves of
\cref{fig:ch11-push-swap}(b) cannot be parallelised into fewer than four
time steps without a conflict, and that the multipush can be parallelised
into $d$ time steps.
\end{exercise}

\begin{exercise}[\difficulty{3} Coding]
\label{exr:ch11-coding}
This is the coding exercise of this chapter in the study plan
(\cref{ch:appA}, Week~5). (a) Implement joint \astar for $k$ agents on a
grid from scratch, with the heuristic \cref{eq:ch11-heuristic} and the
cost \cref{eq:ch11-cost}, and check it against \code{joint\_astar} of the
chapter code on random $5 \times 4$ instances with two to four agents.
(b) Implement independence detection on top of it, with Standley's
refinement that tries to re-plan one group around the other's plan at equal
cost before merging; count how often the refinement avoids a merge on
random $8 \times 8$ maps with six agents, and reproduce the right panel of
\cref{fig:ch11-expansions}. (c) Implement the push and swap primitives on
an abstract graph (adjacency lists, one move at a time), including the
reversal by position, and use them to solve the swap example and the push
example of the chapter; validate every plan with a conflict checker and
compare its makespan and sum of costs with joint \astar. (d) Optional:
implement \code{rotate} and construct an instance with a fully occupied
cycle on which your push-and-swap fails and push-and-rotate succeeds.
\end{exercise}

\begin{exercise}[\difficulty{3}]
\label{exr:ch11-unsolvable}
(a) Give an instance with two empty vertices that has no solution, and
explain in one sentence why no algorithm, complete or not, can solve it.
(b) Give an instance with exactly one empty vertex on a cycle of five
vertices in which the agents must be rotated by one position, and explain
why a swap at a vertex of degree three is impossible while a rotation with a
cleared vertex is not available either; what does this say about the
condition ``at least two empty vertices''? (c) \mstar and \cbs run for ever
on some unsolvable instances and terminate on others. Explain why \mstar
always terminates (\cref{thm:ch11-terminates}) while \cbs needs an extra
argument, and name the instance of (a) as a test case.
\end{exercise}
